\chapter*{Conclusion}
\phantomsection
\addcontentsline{toc}{chapter}{Conclusion}

This thesis introduced the solution for efficient, directly addressable access to the NVMe from the
FPGA accelerator as an alternative to traditional staging of data in the host memory. The
proposed implementation demonstrated shared CPU-FPGA filesystem access to an NVMe device using PCI
Express peer-to-peer communication, with the CPU primarily responsible for configuration and
standard filesystem operation. The core objective was to reduce avoidable software overhead that the
copying of data involves as well as the NVMe command creation and memory-bus pressure while
preserving practical integration with a Linux software stack.

To achieve this, the work combined hardware and software co-design. On the hardware side, the DMA
Iuventus IP core was developed and integrated into the FPGA design to process NVMe-related PCIe
traffic, maintain queue state, and serve requests from accelerator-mapped queue and buffer
regions. On the software side, SPDK with Linux ublk was used to create a userspace control path to
a mountable block device, enabling validation with a standard filesystem workflow instead of only
raw block-level testing. This architecture established an operational path for host-bypassing data
exchange between the FPGA and NVMe controller while preserving host access to shared storage.

The implementation was validated on the AMD-based computer architecture equipped with the \emph{AMD
Alveo U55C} platform and an NVMe SSD. The reported measurements and functional tests confirmed
correct operation of command submission and completion handling, successful read/write behavior for
both the CPU and accelerator access paths, and latency characteristics consistent with expected NVMe
operation. Resource utilization and timing results showed that the proposed hardware is
implementable on a practical data-center accelerator.

At the same time, the solution has limitations emerging from its current scope. The presented setup
was evaluated on a specific hardware and firmware environment. The implementation choices (like
IOMMU settings or only one flying command from the DMA Iuventus module) prioritize functionality and
experimental control over multi-tenant or high throughput use-cases.

Future work can therefore focus on strengthening deployability and scaling. Promising directions
include operation with enabled IOMMU where \verb|vfio| driver can be used to access the unbound
devices for non-superusers, extension toward more parallel queue configurations including queues
with larger size that would allow more flying commands from the hardware, or other models of NVMe
drives. The testbed used in this thesis provided 4 different SSDs but only some of them worked at a
time in the communication with the hardware (only 1 at the time when this thesis was written).
Although all of these NVMe drives were normally operable from the host system, the FPGA access to
all of them was not possible all of the time. The cause of this remains largely undiscovered.
